% ============================================================
%  Informe – Proyecto Final de Cátedra 2026
%  Realidad Virtual – UNCuyo Ing.
% ============================================================
\documentclass[12pt, a4paper]{article}

% --- Encoding & Language ---
\usepackage[utf8]{inputenc}
\usepackage[T1]{fontenc}
\usepackage[spanish, es-tabla]{babel}

% --- Page Layout ---
\usepackage{geometry}
\geometry{
    a4paper,
    top=3.2cm,
    bottom=2.8cm,
    left=2.5cm,
    right=2.5cm,
    headheight=1.6cm,
    headsep=0.6cm,
    footskip=1.2cm
}

% --- Graphics & Color ---
\usepackage{graphicx}
\usepackage[dvipsnames, table]{xcolor}

% --- Header / Footer ---
\usepackage{fancyhdr}
\usepackage{lastpage}

% --- Math & Other Utilities ---
\usepackage{amsmath, amssymb}
\usepackage{float}
\usepackage{url}
\usepackage{booktabs}
\usepackage{caption}

% --- Código fuente (estilo Go) ---
\usepackage{listings}

\definecolor{gogreen}{rgb}{0.0, 0.50, 0.0}
\definecolor{goblue}{rgb}{0.0, 0.0, 0.75}
\definecolor{gogray}{rgb}{0.5, 0.5, 0.5}
\definecolor{gobg}{rgb}{0.97, 0.97, 0.97}

\lstdefinestyle{golang}{
    language=Go,
    backgroundcolor=\color{gobg},
    basicstyle=\ttfamily\footnotesize,
    keywordstyle=\color{goblue}\bfseries,
    commentstyle=\color{gogreen}\itshape,
    stringstyle=\color{red},
    numberstyle=\tiny\color{gogray},
    numbers=left,
    numbersep=6pt,
    frame=single,
    framerule=0.4pt,
    rulecolor=\color{gogray},
    breaklines=true,
    breakatwhitespace=false,
    tabsize=4,
    showstringspaces=false,
    captionpos=b,
}

% \img: incluye imagen (el archivo debe existir)
\newcommand{\img}[2][width=\textwidth]{%
    \includegraphics[#1]{#2}%
}

% \imgpending: placeholder para imágenes aún no disponibles
\newcommand{\imgpending}[1]{%
    \begin{center}
        \fbox{\begin{minipage}{0.92\linewidth}
            \vspace{1.2cm}
            \centering\itshape\small
            [Imagen pendiente: \texttt{#1}]
            \vspace{1.2cm}
        \end{minipage}}
    \end{center}%
}

\usepackage{hyperref}
\hypersetup{hidelinks}

% ============================================================
%  COLORES
% ============================================================
\definecolor{azulUNC}{RGB}{0, 70, 127}
\definecolor{verdeRV}{RGB}{0, 110, 100}
\definecolor{grisLinea}{RGB}{90, 90, 90}

% ============================================================
%  ENCABEZADO Y PIE  (páginas del cuerpo)
% ============================================================
\pagestyle{fancy}
\fancyhf{}

\fancyhead[L]{%
    \footnotesize
    UNCuyo -- Facultad de Ingeniería \\
    Mendoza -- Argentina
}

\fancyhead[C]{%
    \footnotesize\bfseries
    Realidad Virtual (00318) \\
    PROYECTO FINAL DE CÁTEDRA
}

\fancyhead[R]{%
    \footnotesize
    Escobar, M. -- Martín Duci, I. \\
    Julio 2026
}

\fancyfoot[C]{%
    \small\bfseries Pág.\ \thepage\ de \pageref{LastPage}
}

\renewcommand{\headrulewidth}{0.6pt}
\renewcommand{\footrulewidth}{0.6pt}

\renewcommand{\headrule}{%
    {\color{grisLinea}\hrule width\headwidth height\headrulewidth}%
    \vskip -\headrulewidth
}
\renewcommand{\footrule}{%
    {\color{grisLinea}\hrule width\headwidth height\footrulewidth}%
    \vskip -\footrulewidth
}

% ============================================================
%  DOCUMENTO
% ============================================================
\begin{document}

% ── PORTADA ─────────────────────────────────────────────────
\thispagestyle{empty}

\begin{center}

    \vspace*{1.5cm}

    {\large\bfseries
        Desarrollo de una Arquitectura Cliente-Servidor HTTP\\[5pt]
        para el Cálculo y Simulación de Ecuaciones Diferenciales\\[5pt]
        mediante Realidad Virtual
    }

    \vspace{1.8cm}

    {\LARGE\bfseries\color{azulUNC} Universidad Nacional de Cuyo}\\[6pt]
    {\large Facultad de Ingeniería}

    \vspace{3cm}

    {\large\color{verdeRV}\bfseries Realidad Virtual (00318)}

    \vspace{0.4cm}

    {\normalsize Proyecto Final de Cátedra}

    \vspace{3cm}

    {\normalsize\bfseries Profesores}\\[6pt]
    {\normalsize\color{azulUNC}
        Prof. Rosenstein, Javier \\[4pt]
        Prof. Tomba, Carlos
    }

    \vspace{1.2cm}

    {\normalsize\bfseries Alumnos}\\[6pt]
    {\normalsize\color{azulUNC}
        Escobar, Matías \quad -- \quad XXXXX \\[4pt]
        Martín Duci, Ignacio \quad -- \quad 13560
    }

    \vspace{1.2cm}

    {\normalsize Julio 2026}

\end{center}

% ── ÍNDICE ──────────────────────────────────────────────────
\newpage
\thispagestyle{fancy}
\tableofcontents

% ── RESUMEN ─────────────────────────────────────────────────
\newpage
\thispagestyle{fancy}

\begin{center}
    {\large\bfseries Resumen}
\end{center}

\vspace{0.6cm}

El presente trabajo describe el desarrollo de un sistema cliente-servidor orientado al
entorno académico para el cálculo numérico y la visualización interactiva de ecuaciones
diferenciales parciales (EDPs). El sistema fue concebido como una plataforma de
exploración y aprendizaje que permite a los alumnos simular, probar e iterar sobre
distintas configuraciones de EDPs — variando parámetros físicos, condiciones iniciales
y dominios espaciales — para construir intuición sobre su comportamiento de forma
activa. Como funcionalidad complementaria, el sistema incorpora un mecanismo de
sincronización que permite al profesor propagar una configuración específica a todos
los dispositivos conectados, facilitando la discusión guiada en clase a partir de un
caso común.

El proyecto se estructura en dos etapas diferenciadas. La primera corresponde al
\textbf{servidor}, implementado en el lenguaje \textbf{Go} utilizando el framework
\textbf{Gin} para la gestión de rutas HTTP. El servidor expone una API REST que permite
el registro de usuarios, el cálculo de soluciones numéricas y el control de
sincronización entre dispositivos. Se implementaron dos solvers de ecuaciones
diferenciales parciales: la \textbf{ecuación de onda 2D}, resuelta mediante el esquema
explícito \textit{leapfrog} de segundo orden con condición CFL, y la
\textbf{ecuación de calor 2D} sobre una placa cuadrada, resuelta mediante el esquema
\textit{FTCS} (\textit{Forward Time Centered Space}) explícito. Ambos solvers operan
sobre un dominio espacial discretizado con condiciones de contorno de Dirichlet, y
emplean técnicas de \textit{rolling buffer} y submuestreo temporal para mantener el
uso de memoria acotado independientemente de la duración de la simulación. El servidor
incluye además un panel web docente desarrollado en HTML, CSS y JavaScript, que permite
visualizar las soluciones en una representación tridimensional renderizada por software
mediante un rasterizador de polígonos implementado sobre el elemento \texttt{Canvas} de
HTML5, con controles de órbita y un sistema de sincronización de parámetros hacia los
alumnos.

La segunda etapa corresponde a la \textbf{aplicación móvil de realidad virtual},
desarrollada en \textbf{Unity}. La aplicación consume la API del servidor y presenta
las soluciones numéricas en un entorno de realidad virtual inmersivo, permitiendo al
alumno interactuar con la simulación en tiempo real. La visualización en realidad
virtual aporta una dimensión perceptual adicional respecto a las representaciones
convencionales, facilitando la comprensión intuitiva del comportamiento de las EDPs en
el espacio y el tiempo.

\vfill

% ── SECCIÓN 1: INTRODUCCIÓN ─────────────────────────────────
\newpage

\section{Introducción}

Las ecuaciones diferenciales parciales (EDPs) constituyen una de las herramientas
matemáticas fundamentales en ingeniería y ciencias aplicadas. Modelan fenómenos
físicos tan diversos como la propagación de ondas, la difusión de calor, el flujo de
fluidos y la deformación de estructuras. Sin embargo, su comprensión intuitiva
representa un desafío significativo para los estudiantes, en parte debido a que los
métodos de resolución analítica se limitan a casos muy particulares y las soluciones
numéricas suelen presentarse de forma estática en hojas de cálculo o gráficas
bidimensionales.

El presente proyecto propone un enfoque alternativo: una plataforma interactiva que
permite al usuario simular el comportamiento de EDPs en tiempo real, variar libremente
sus parámetros y condiciones iniciales, y observar las soluciones mediante
visualizaciones tridimensionales inmersivas en realidad virtual. El objetivo no es
reemplazar la teoría sino complementarla, ofreciendo un entorno donde la experimentación
directa genere intuición física sobre fenómenos que de otro modo resultan abstractos.

El sistema se compone de dos componentes principales que se comunican mediante una
arquitectura cliente-servidor HTTP. El servidor, desarrollado en Go, se encarga del
cálculo numérico de las soluciones y de gestionar múltiples usuarios simultáneos con
sus respectivas sesiones. El cliente, una aplicación móvil de realidad virtual
desarrollada en Unity, consume las soluciones del servidor y las presenta al usuario en
un entorno inmersivo e interactivo.

Como funcionalidad adicional orientada al uso en aula, el sistema incorpora un
mecanismo de sincronización que permite al profesor propagar una configuración
específica a todos los dispositivos conectados, habilitando la discusión guiada sobre
un caso común.

% ── SECCIÓN 2: OBJETIVOS Y BENEFICIARIOS ────────────────────
\newpage

\section{Objetivos y Beneficiarios}

\subsection{Objetivo general}

Desarrollar una plataforma cliente-servidor para la simulación interactiva de
ecuaciones diferenciales parciales, con visualización en realidad virtual, que
facilite el aprendizaje activo mediante la exploración paramétrica de soluciones
numéricas en tiempo real.

\subsection{Objetivos específicos}

\begin{itemize}
    \item Implementar solvers numéricos estables y eficientes para la ecuación de onda
    2D y la ecuación de calor 2D sobre dominios cuadrados.

    \item Diseñar una API HTTP que exponga los resultados de la simulación a clientes
    heterogéneos (app móvil, panel web) de forma eficiente y con control de acceso por
    roles.

    \item Desarrollar una aplicación móvil en Unity que presente las soluciones
    numéricas en un entorno de realidad virtual inmersivo, permitiendo al usuario
    interactuar con la simulación.

    \item Proveer un panel web docente que permita visualizar, controlar y sincronizar
    simulaciones hacia múltiples usuarios conectados.

    \item Garantizar la estabilidad numérica de los esquemas implementados mediante el
    cumplimiento de las condiciones de estabilidad correspondientes a cada método.
\end{itemize}

\subsection{Beneficiarios}

La plataforma está orientada principalmente a \textbf{estudiantes de carreras de
ingeniería y ciencias exactas} que cursen asignaturas donde las EDPs tienen un rol
central: mecánica de medios continuos, transferencia de calor, electromagnetismo,
mecánica de fluidos, entre otras. La posibilidad de explorar libremente parámetros y
condiciones iniciales convierte al sistema en una herramienta de estudio autónomo que
complementa la clase magistral.

En segundo lugar, la plataforma beneficia a los \textbf{docentes}, que pueden utilizarla
como recurso didáctico durante la clase para ilustrar conceptos con simulaciones en
vivo, o proponer a los alumnos escenarios específicos mediante la funcionalidad de
sincronización.

% ── SECCIÓN 3: RESOLUCIÓN DE EDPs ───────────────────────────
\newpage

\section{Resolución de Ecuaciones Diferenciales}

\subsection{Fundamentos del método de diferencias finitas}

Las EDPs implementadas en este proyecto se resuelven mediante el \textbf{método de
diferencias finitas} (MDF), que consiste en discretizar el dominio continuo en una
grilla de puntos y aproximar las derivadas parciales mediante cocientes de diferencias
entre valores en nodos vecinos.

Dado un dominio espacial cuadrado $[0,L]\times[0,L]$ con $n$ subdivisiones, el paso
espacial es $\Delta x = \Delta y = L/n$. Los nodos interiores son los puntos
$(i\Delta x,\, j\Delta y)$ con $1 \leq i,j \leq n-2$. Las condiciones de contorno de
Dirichlet se imponen directamente sobre los nodos del borde, fijando su valor para
todo tiempo.

La aproximación por diferencias centradas de segundo orden para la derivada segunda
espacial en 2D es el \textbf{laplaciano discreto}:

\begin{equation}
    \nabla^2 u_{i,j} \approx
    \frac{u_{i+1,j} + u_{i-1,j} + u_{i,j+1} + u_{i,j-1} - 4\,u_{i,j}}{\Delta x^2}
    \label{eq:laplaciano}
\end{equation}

Esta expresión es la base de ambos esquemas numéricos implementados.

% ---- 3.1 -------------------------------------------------------
\subsection{Ecuación de onda 2D}

\subsubsection{Formulación matemática}

La ecuación de onda 2D describe la evolución temporal de una membrana vibrante:

\begin{equation}
    \frac{\partial^2 u}{\partial t^2} = c^2
    \left(\frac{\partial^2 u}{\partial x^2} + \frac{\partial^2 u}{\partial y^2}\right)
    \label{eq:onda2d}
\end{equation}

donde $u(x,y,t)$ es el desplazamiento de la membrana y $c$ es la velocidad de
propagación de la onda. Se imponen condiciones de contorno de Dirichlet homogéneas
($u = 0$ en todo el borde) y velocidad inicial nula ($\partial u/\partial t\big|_{t=0}
= 0$).

\subsubsection{Esquema leapfrog de segundo orden}

Se utiliza el esquema explícito \textit{leapfrog} centrado en tiempo y espacio, que
aproxima la derivada segunda temporal:

\begin{equation}
    \frac{\partial^2 u}{\partial t^2} \approx
    \frac{u^{k+1}_{i,j} - 2\,u^k_{i,j} + u^{k-1}_{i,j}}{\Delta t^2}
\end{equation}

Combinando con el laplaciano discreto~\eqref{eq:laplaciano} y despejando $u^{k+1}$:

\begin{equation}
    u^{k+1}_{i,j} = 2\,u^k_{i,j} - u^{k-1}_{i,j}
    + r^2\!\left(u^k_{i+1,j} + u^k_{i-1,j} + u^k_{i,j+1} + u^k_{i,j-1}
    - 4\,u^k_{i,j}\right)
    \label{eq:leapfrog}
\end{equation}

donde $r = c\,\Delta t / \Delta x$ es el \textbf{número de Courant}: mide cuántas
celdas espaciales recorre la perturbación en un paso temporal. Un valor $r > 1/\sqrt{2}$
significa que la perturbación ``salta'' más rápido de lo que el esquema puede
representar, provocando inestabilidad numérica.

Para el primer paso temporal, donde no existe $u^{k-1}$, se aplica la condición de
velocidad inicial nula para obtener un arranque de primer orden:

\begin{equation}
    u^1_{i,j} = u^0_{i,j} + \frac{r^2}{2}
    \left(u^0_{i+1,j} + u^0_{i-1,j} + u^0_{i,j+1} + u^0_{i,j-1}
    - 4\,u^0_{i,j}\right)
\end{equation}

\subsubsection{Condición de estabilidad CFL}

Para que el esquema leapfrog sea estable en 2D, el número de Courant debe satisfacer:

\begin{equation}
    r = \frac{c\,\Delta t}{\Delta x} \leq \frac{1}{\sqrt{2}}
    \label{eq:cfl-onda}
\end{equation}

En la implementación se usa $r = 0.9/\sqrt{2}$ para mantener un margen de seguridad.
El paso temporal resultante es $\Delta t = r\,\Delta x / c$, y el número de pasos
temporales necesario para simular $T$ segundos es $n_t = \lceil T / \Delta t \rceil$.

\subsubsection{Condiciones iniciales implementadas}

Se ofrecen tres configuraciones de condición inicial $u(x,y,0)$:

\begin{itemize}
    \item \textbf{Seno}: modo fundamental de vibración de la membrana cuadrada:
    \[u_0 = \sin\!\left(\frac{\pi x}{L}\right)\sin\!\left(\frac{\pi y}{L}\right)\]
    \item \textbf{Gaussiana}: pulso centrado que se propaga como ondas circulares:
    \[u_0 = \exp\!\left[-\left(\frac{x-L/2}{L/5}\right)^{\!2}
            -\left(\frac{y-L/2}{L/5}\right)^{\!2}\right]\]
    \item \textbf{Triangular}: pirámide lineal centrada:
    \[u_0 = \bigl(1-|2x/L-1|\bigr)\bigl(1-|2y/L-1|\bigr)\]
\end{itemize}

\subsubsection{Implementación en Go}

\begin{lstlisting}[style=golang, caption={Solver de la ecuación de onda 2D — \texttt{solver.go}}, label=lst:onda2d]
func ecuacion_onda_2d(L, T, c float64, inicial string) [][][]float64 {
    n := int(math.Ceil(L * 8))
    if n < 16 { n = 16 }
    d := L / float64(n)

    CFL := 0.9 / math.Sqrt2
    d_t := CFL * d / c
    n_t := int(math.Ceil(T / d_t))
    d_t  = T / float64(n_t)
    r   := (c * d_t / d) * (c * d_t / d)

    maxFrames := 200
    step := 1
    if n_t > maxFrames { step = n_t / maxFrames }

    out := make([][][]float64, 0, maxFrames)
    pp  := newGrid(n)  // t-2
    p   := newGrid(n)  // t-1
    cur := newGrid(n)  // t

    // Condicion inicial
    for i := 1; i < n-1; i++ {
        for j := 1; j < n-1; j++ {
            x := float64(i) * d
            y := float64(j) * d
            switch inicial {
            case "gauss":
                pp[i][j] = math.Exp(-math.Pow((x-L/2)/(L/5), 2) -
                                     math.Pow((y-L/2)/(L/5), 2))
            case "triangular":
                tx := 1 - math.Abs(2*x/L-1)
                ty := 1 - math.Abs(2*y/L-1)
                pp[i][j] = tx * ty
            default: // seno
                pp[i][j] = math.Sin(math.Pi*x/L) * math.Sin(math.Pi*y/L)
            }
        }
    }
    out = append(out, copyGrid(pp))

    // Primer paso (velocidad inicial = 0)
    for i := 1; i < n-1; i++ {
        for j := 1; j < n-1; j++ {
            p[i][j] = pp[i][j] + 0.5*r*(pp[i+1][j]+pp[i-1][j]+
                       pp[i][j+1]+pp[i][j-1]-4*pp[i][j])
        }
    }
    if step == 1 { out = append(out, copyGrid(p)) }

    // Pasos restantes - leapfrog
    for t := 2; t < n_t; t++ {
        for i := 1; i < n-1; i++ {
            for j := 1; j < n-1; j++ {
                cur[i][j] = 2*p[i][j] - pp[i][j] +
                    r*(p[i+1][j]+p[i-1][j]+p[i][j+1]+p[i][j-1]-4*p[i][j])
            }
        }
        if t%step == 0 { out = append(out, copyGrid(cur)) }
        pp, p, cur = p, cur, pp
    }
    return out
}
\end{lstlisting}

% ---- 3.2 -------------------------------------------------------
\subsection{Ecuación de calor 2D}

\subsubsection{Formulación matemática}

La ecuación de calor 2D describe la difusión de temperatura sobre una placa cuadrada:

\begin{equation}
    \frac{\partial u}{\partial t} = \alpha
    \left(\frac{\partial^2 u}{\partial x^2} + \frac{\partial^2 u}{\partial y^2}\right)
    \label{eq:calor2d}
\end{equation}

donde $u(x,y,t)$ es la temperatura normalizada y $\alpha$ es la difusividad térmica del
material.

\subsubsection{Esquema FTCS explícito}

Se utiliza el esquema \textit{Forward Time Centered Space} (FTCS), que aproxima la
derivada temporal mediante diferencias hacia adelante de primer orden:

\begin{equation}
    \frac{\partial u}{\partial t} \approx
    \frac{u^{k+1}_{i,j} - u^k_{i,j}}{\Delta t}
\end{equation}

Combinando con el laplaciano discreto~\eqref{eq:laplaciano} y despejando $u^{k+1}$:

\begin{equation}
    u^{k+1}_{i,j} = u^k_{i,j} + r_\alpha\!
    \left(u^k_{i+1,j} + u^k_{i-1,j} + u^k_{i,j+1} + u^k_{i,j-1}
    - 4\,u^k_{i,j}\right)
    \label{eq:ftcs}
\end{equation}

donde $r_\alpha = \alpha\,\Delta t / \Delta x^2$ es el número de difusión.

A diferencia del esquema leapfrog, el FTCS solo requiere el instante anterior para
avanzar en tiempo, lo que simplifica la implementación al necesitar únicamente dos
grillas en lugar de tres.

\subsubsection{Condición de estabilidad}

Para que el esquema FTCS sea estable en 2D, el número de difusión debe satisfacer:

\begin{equation}
    r_\alpha = \frac{\alpha\,\Delta t}{\Delta x^2} \leq \frac{1}{4}
    \label{eq:estab-calor}
\end{equation}

En la implementación se usa $r_\alpha = 0.20$ como margen de seguridad respecto al
límite teórico de $0.25$.

\subsubsection{Condiciones iniciales y de contorno implementadas}

Se ofrecen dos configuraciones:

\begin{itemize}
    \item \textbf{Gaussiana}: pulso de temperatura centrado. Los cuatro bordes se fijan
    en $u = 0$ (sumideros fríos). El calor se disipa hacia los bordes y la temperatura
    interior decae hasta equilibrarse en cero.

    \item \textbf{Bordes calientes}: los cuatro bordes se fijan en $u = 1$ (fuentes de
    calor constantes) y el interior comienza en $u = 0$. El calor penetra desde las
    paredes hacia el centro hasta alcanzar el estado estacionario $u = 1$ en todo el
    dominio.
\end{itemize}

En ambos casos las condiciones de contorno se reimponen explícitamente después de cada
paso temporal para mantener el valor fijo en el borde independientemente de la difusión
numérica.

\subsubsection{Implementación en Go}

\begin{lstlisting}[style=golang, caption={Solver de la ecuación de calor 2D — \texttt{solver.go}}, label=lst:calor2d]
func ecuacion_calor_2d(L, T, alpha float64, inicial string) [][][]float64 {
    n := int(math.Ceil(L * 5))
    if n < 10 { n = 10 }
    d := L / float64(n)

    r_max := 0.20
    d_t   := r_max * d * d / alpha
    n_t   := int(math.Ceil(T / d_t))
    d_t    = T / float64(n_t)
    r     := alpha * d_t / (d * d)

    maxFrames := 200
    step := 1
    if n_t > maxFrames { step = n_t / maxFrames }

    out  := make([][][]float64, 0, maxFrames)
    prev := newGrid(n)
    curr := newGrid(n)

    bcVal := 0.0
    if inicial == "bordes" { bcVal = 1.0 }

    // Condicion inicial
    for i := 1; i < n-1; i++ {
        for j := 1; j < n-1; j++ {
            x := float64(i) * d
            y := float64(j) * d
            switch inicial {
            case "gauss":
                prev[i][j] = math.Exp(-(math.Pow((x-L/2)/(L/5), 2) +
                                         math.Pow((y-L/2)/(L/5), 2)))
            default: // bordes calientes: interior frio
                prev[i][j] = 0.0
            }
        }
    }
    // Aplicar condicion de contorno inicial
    for k := 0; k < n; k++ {
        prev[k][0] = bcVal; prev[k][n-1] = bcVal
        prev[0][k] = bcVal; prev[n-1][k] = bcVal
    }
    out = append(out, copyGrid(prev))

    // Avance temporal FTCS
    for t := 1; t < n_t; t++ {
        for i := 1; i < n-1; i++ {
            for j := 1; j < n-1; j++ {
                curr[i][j] = prev[i][j] + r*(prev[i+1][j]+prev[i-1][j]+
                              prev[i][j+1]+prev[i][j-1]-4*prev[i][j])
            }
        }
        // Reimposicion de condicion de contorno
        for k := 0; k < n; k++ {
            curr[k][0] = bcVal; curr[k][n-1] = bcVal
            curr[0][k] = bcVal; curr[n-1][k] = bcVal
        }
        if t%step == 0 { out = append(out, copyGrid(curr)) }
        prev, curr = curr, prev
    }
    return out
}
\end{lstlisting}

\subsection{Consideraciones de eficiencia}

Ambos solvers comparten dos estrategias para mantener el uso de recursos acotado
independientemente de la duración o resolución de la simulación:

\begin{itemize}
    \item \textbf{Rolling buffer}: en lugar de almacenar todos los pasos temporales en
    memoria, se mantienen únicamente dos o tres grillas activas que se rotan mediante
    reasignación de punteros (\texttt{pp, p, cur = p, cur, pp}). El costo de memoria
    es $O(n^2)$ constante.

    \item \textbf{Submuestreo temporal}: si la simulación requiere más de 200 pasos
    temporales, solo se guarda uno de cada $\lfloor n_t / 200 \rfloor$ pasos. Esto
    limita el tamaño de la respuesta JSON a un máximo de 200 frames sin afectar la
    estabilidad numérica, que depende del paso $\Delta t$ interno y no de cuántos
    frames se exportan.
\end{itemize}

% ── SECCIÓN 4: ARQUITECTURA DEL SERVIDOR ────────────────────
\newpage

\section{Arquitectura del Servidor}

El servidor está escrito íntegramente en Go y utiliza el framework \textbf{Gin} para la
gestión de rutas HTTP. El código se organiza en cinco archivos, cada uno con una
responsabilidad bien delimitada:

\begin{center}
\begin{tabular}{ll}
\toprule
\textbf{Archivo} & \textbf{Responsabilidad} \\
\midrule
\texttt{main.go}     & Punto de entrada: router y arranque del servidor \\
\texttt{auth.go}     & Autenticación del profesor y limitación de tasa \\
\texttt{sync.go}     & Estado global en memoria: clientes, roles y sincronización \\
\texttt{handlers.go} & Endpoints de la app móvil \\
\texttt{admin.go}    & Endpoints del panel web docente \\
\texttt{solver.go}   & Solvers numéricos de EDPs (sin dependencias HTTP) \\
\bottomrule
\end{tabular}
\end{center}

Todos los archivos pertenecen al mismo paquete \texttt{main}, lo que significa que
comparten tipos, variables y funciones sin necesidad de importaciones entre ellos. El
estado global del servidor vive en una única variable (\texttt{estado}) definida en
\texttt{sync.go}, y todos los handlers acceden a ella mediante sus métodos.

El diagrama siguiente resume cómo se conectan los archivos:

\begin{verbatim}
                 +-------------+
                 |   main.go   |  define rutas, arranca :8080
                 +------+------+
       +-----------------+------------------+
       |                 |                  |
  +---------+    +------------+    +----------+
  | auth.go |    | handlers.go|    | admin.go |
  |  /auth  |    | /register  |    | /admin/* |
  |  /panel |    | /solve     |    +----+-----+
  +---------+    | /sync/*    |         |
                 +------+-----+         |
                        |               |
                 +------+---------------+--+
                 |         sync.go         |
                 |  estado global (RWMu)   |
                 +------------+------------+
                              |
                       +------+------+
                       |  solver.go  |
                       | onda/calor  |
                       +-------------+
\end{verbatim}

\subsection{Resumen de rutas}

La tabla siguiente sintetiza todos los endpoints expuestos por el servidor, el método
HTTP, quién puede usarlos y qué devuelven.

\begin{center}
\renewcommand{\arraystretch}{1.25}
\begin{tabular}{p{3.6cm} p{1.4cm} p{3.2cm} p{5.2cm}}
\toprule
\textbf{Ruta} & \textbf{Método} & \textbf{Acceso} & \textbf{Función / Respuesta} \\
\midrule
\texttt{/}
  & GET & Todos
  & Sirve la pantalla de login del profesor (\texttt{login.html}) \\
\texttt{/auth}
  & POST & Todos
  & Verifica el código del profesor y autoriza su IP. Devuelve \texttt{\{ok:true\}} \\
\texttt{/panel}
  & GET & Profesor (IP autorizada)
  & Sirve el panel docente (\texttt{panel.html}). Redirige al login si la IP no está autorizada \\
\texttt{/register}
  & POST & App móvil
  & Registra o actualiza un alumno (legajo, nombre). Devuelve \texttt{\{ok:true\}} \\
\texttt{/solve}
  & GET & App móvil (1 req/5s)
  & Calcula la solución de la EDP con los parámetros dados. Devuelve \texttt{\{sync, ecuacion, result\}} con hasta 200 frames \\
\texttt{/sync/set}
  & POST & Asistente
  & Activa la sincronización con los parámetros indicados. Solo accesible con rol de asistente \\
\texttt{/sync/clear}
  & POST & Asistente
  & Desactiva la sincronización. Solo accesible con rol de asistente \\
\midrule
\texttt{/admin/clients}
  & GET & Profesor
  & Lista todos los alumnos registrados con IP, legajo, nombre, rol y hora del último request \\
\texttt{/admin/promote}
  & POST & Profesor
  & Otorga rol de asistente al cliente con la IP indicada \\
\texttt{/admin/demote}
  & POST & Profesor
  & Revoca el rol de asistente del cliente con la IP indicada \\
\texttt{/admin/sync/set}
  & POST & Profesor
  & Activa la sincronización directamente desde el panel, sin requerir asistente \\
\texttt{/admin/sync/clear}
  & POST & Profesor
  & Desactiva la sincronización desde el panel \\
\bottomrule
\end{tabular}
\end{center}

% -------- main.go -----------------------------------------------
\subsection{\texttt{main.go} — Router y punto de entrada}

\texttt{main.go} es el punto de entrada del servidor. Declara la constante del código
de acceso del profesor y en \texttt{main()} crea el router, registra todas las rutas y
arranca el servidor en el puerto 8080.

\begin{lstlisting}[style=golang, caption={\texttt{main.go} completo}]
package main

import "github.com/gin-gonic/gin"

const codigoProfesor = "123"

func main() {
    router := gin.Default()
\end{lstlisting}

\texttt{gin.Default()} crea un router con los middlewares de registro de requests y
recuperación de panics ya activados. El código del profesor está hardcodeado por
simplicidad: en el contexto de aula esto es aceptable, ya que la red es controlada y
el código se cambia directamente en el binario si es necesario.

\begin{lstlisting}[style=golang]
    router.StaticFile("/", "./web/login.html")
    router.POST("/auth",  authHandler)
    router.GET("/panel",  panelHandler)
\end{lstlisting}

Estas tres rutas forman el flujo de autenticación del profesor. La raíz sirve la
pantalla de login. \texttt{POST /auth} valida el código y registra la IP. \texttt{GET
/panel} sirve el panel solo si la IP ya fue autorizada.

\begin{lstlisting}[style=golang]
    admin := router.Group("/admin")
    {
        admin.GET("/clients",     clientesHandler)
        admin.POST("/promote",    promoverHandler)
        admin.POST("/demote",     demoverHandler)
        admin.POST("/sync/clear", syncClearHandler)
        admin.POST("/sync/set",   syncSetAdminHandler)
    }
\end{lstlisting}

El grupo \texttt{/admin} agrupa todas las rutas exclusivas del panel docente bajo un
prefijo común. Si en el futuro se quisiera añadir un middleware de verificación de IP
a todas ellas, bastaría con agregarlo al grupo sin tocar cada handler.

\begin{lstlisting}[style=golang]
    router.POST("/register", registerHandler)
    router.GET("/solve", RateLimitMiddleware(), solveHandler)
    router.POST("/sync/set",   syncSetHandler)
    router.POST("/sync/clear", syncClearAsistenteHandler)

    router.Run("10.83.158.17:8080")
}
\end{lstlisting}

Las rutas de la app móvil son públicas: cualquier dispositivo en la red puede
acceder a ellas. La excepción es \texttt{/solve}, que pasa primero por
\texttt{RateLimitMiddleware()}: este middleware limita a una solicitud cada cinco
segundos por IP, evitando que un alumno sature el servidor con cálculos costosos
consecutivos.

% -------- auth.go -----------------------------------------------
\subsection{\texttt{auth.go} — Autenticación y limitación de tasa}

Este archivo implementa dos mecanismos de control de acceso independientes: la
autenticación del profesor por IP y el limitador de tasa para el endpoint de cálculo.

\begin{lstlisting}[style=golang, caption={\texttt{auth.go} — imports y variables globales}]
package main

import (
    "fmt"
    "math"
    "net/http"
    "sync"
    "time"

    "github.com/gin-gonic/gin"
    "golang.org/x/time/rate"
)

type authRequest struct {
    Codigo string `json:"codigo"`
}

var (
    ipsAutorizadas   = make(map[string]bool)
    ipsAutorizadasMu sync.RWMutex
)

var limiters sync.Map
\end{lstlisting}

\texttt{ipsAutorizadas} es el mapa que actúa como lista blanca de IPs del profesor.
Está protegido por un \texttt{sync.RWMutex} propio (distinto del mutex del estado
global) porque es independiente del ciclo de vida de los clientes.
\texttt{limiters} es un \texttt{sync.Map} — un mapa diseñado para acceso concurrente
sin necesidad de mutex manual — que almacena un limitador de tasa por IP.

\begin{lstlisting}[style=golang, caption={Limitador de tasa por IP}]
func getLimiter(ip string) *rate.Limiter {
    if v, ok := limiters.Load(ip); ok {
        return v.(*rate.Limiter)
    }
    lim := rate.NewLimiter(rate.Every(5*time.Second), 1)
    actual, _ := limiters.LoadOrStore(ip, lim)
    return actual.(*rate.Limiter)
}

func RateLimitMiddleware() gin.HandlerFunc {
    return func(c *gin.Context) {
        ip := c.ClientIP()
        lim := getLimiter(ip)

        res := lim.Reserve()
        if !res.OK() {
            c.AbortWithStatusJSON(http.StatusTooManyRequests,
                gin.H{"error": "Demasiadas peticiones"})
            return
        }

        delay := res.Delay()
        if delay <= 0 {
            c.Next()
            return
        }

        res.Cancel()
        sec := int(math.Ceil(delay.Seconds()))
        c.Header("Retry-After", fmt.Sprintf("%d", sec))
        c.AbortWithStatusJSON(http.StatusTooManyRequests, gin.H{
            "error":       "Demasiadas peticiones",
            "waitSeconds": sec,
            "message":     fmt.Sprintf("Espera %ds antes de reintentar", sec),
        })
    }
}
\end{lstlisting}

\texttt{rate.NewLimiter(rate.Every(5*time.Second), 1)} crea un limitador de tipo
\textit{token bucket} con capacidad de un token que se regenera cada cinco segundos.
\texttt{Reserve()} consume el token si está disponible y devuelve cuánto tiempo habría
que esperar de lo contrario. Si el delay es mayor que cero, el middleware cancela la
reserva (para no consumir el token inútilmente), agrega la cabecera \texttt{Retry-After}
y rechaza el request con código 429.

\begin{lstlisting}[style=golang, caption={Autenticación del profesor}]
func autorizarIP(ip string) {
    ipsAutorizadasMu.Lock()
    defer ipsAutorizadasMu.Unlock()
    ipsAutorizadas[ip] = true
    fmt.Println("[auth] IP autorizada:", ip)
}

func ipAutorizada(ip string) bool {
    ipsAutorizadasMu.RLock()
    defer ipsAutorizadasMu.RUnlock()
    return ipsAutorizadas[ip]
}

func authHandler(c *gin.Context) {
    var req authRequest
    if err := c.ShouldBindJSON(&req); err != nil {
        c.JSON(http.StatusBadRequest,
            gin.H{"error": "Cuerpo de solicitud invalido"})
        return
    }
    if req.Codigo != codigoProfesor {
        c.JSON(http.StatusUnauthorized, gin.H{"error": "Codigo incorrecto"})
        return
    }
    autorizarIP(c.ClientIP())
    c.JSON(http.StatusOK, gin.H{"ok": true})
}

func panelHandler(c *gin.Context) {
    if !ipAutorizada(c.ClientIP()) {
        c.Redirect(http.StatusFound, "/")
        return
    }
    c.File("./web/panel.html")
}
\end{lstlisting}

\texttt{authHandler} valida el código y delega en \texttt{autorizarIP}, que escribe en
el mapa con lock exclusivo. \texttt{panelHandler} usa \texttt{RLock} (lectura
compartida) para verificar la IP, ya que múltiples profesores podrían abrir el panel
simultáneamente sin interferirse. Si la IP no está autorizada, redirige al login en
lugar de devolver un error, para que el flujo sea transparente en el navegador.

% -------- sync.go -----------------------------------------------
\subsection{\texttt{sync.go} — Estado global en memoria}

Este archivo es el núcleo del servidor: define todos los tipos de datos y centraliza el
estado compartido entre todos los handlers mediante una única variable global
\texttt{estado}.

\begin{lstlisting}[style=golang, caption={\texttt{sync.go} — tipos y estado global}]
package main

import (
    "sync"
    "time"
)

type rol int

const (
    rolAlumno    rol = iota
    rolAsistente rol = iota
)

type Cliente struct {
    IP            string
    Legajo        string
    Nombre        string
    Rol           rol
    UltimoRequest time.Time
}

type parametrosSync struct {
    Ecuacion string
    L        float64
    T        float64
    C        float64
    Inicial  string
}

type estadoGlobal struct {
    mu       sync.RWMutex
    clientes map[string]*Cliente
    sync     struct {
        activo     bool
        parametros parametrosSync
    }
}

var estado = &estadoGlobal{
    clientes: make(map[string]*Cliente),
}
\end{lstlisting}

\texttt{rol} usa el patrón \texttt{iota} de Go para definir constantes enteras
enumeradas. \texttt{estadoGlobal} agrupa clientes y estado de sincronización bajo un
único \texttt{sync.RWMutex}: esto garantiza que todas las operaciones sobre el estado
sean atómicas y que múltiples goroutines (una por request) puedan leer
simultáneamente sin bloquearse entre sí, reservando la exclusividad solo para
escrituras.

\begin{lstlisting}[style=golang, caption={Registro de clientes}]
func (e *estadoGlobal) RegistrarCliente(ip, legajo, nombre string) {
    e.mu.Lock()
    defer e.mu.Unlock()

    if c, existe := e.clientes[ip]; existe {
        c.Legajo = legajo
        c.Nombre = nombre
        c.UltimoRequest = time.Now()
        return
    }

    e.clientes[ip] = &Cliente{
        IP:            ip,
        Legajo:        legajo,
        Nombre:        nombre,
        Rol:           rolAlumno,
        UltimoRequest: time.Now(),
    }
}

func (e *estadoGlobal) ActualizarRequest(ip string) {
    e.mu.Lock()
    defer e.mu.Unlock()
    if c, existe := e.clientes[ip]; existe {
        c.UltimoRequest = time.Now()
    }
}
\end{lstlisting}

\texttt{RegistrarCliente} es idempotente: si la IP ya existe, actualiza legajo y nombre
pero \emph{no} toca el campo \texttt{Rol}. Esto es fundamental para que un asistente
pueda cerrar y reabrir la app sin perder el rol que le otorgó el profesor.

\begin{lstlisting}[style=golang, caption={Consulta y modificacion de clientes}]
func (e *estadoGlobal) Clientes() []*Cliente {
    e.mu.RLock()
    defer e.mu.RUnlock()
    lista := make([]*Cliente, 0, len(e.clientes))
    for _, c := range e.clientes {
        lista = append(lista, c)
    }
    return lista
}

func (e *estadoGlobal) BuscarCliente(ip string) *Cliente {
    e.mu.RLock()
    defer e.mu.RUnlock()
    return e.clientes[ip]
}

func (e *estadoGlobal) PromoverAsistente(ip string) bool {
    e.mu.Lock()
    defer e.mu.Unlock()
    c, existe := e.clientes[ip]
    if !existe { return false }
    c.Rol = rolAsistente
    return true
}

func (e *estadoGlobal) DemoverAsistente(ip string) bool {
    e.mu.Lock()
    defer e.mu.Unlock()
    c, existe := e.clientes[ip]
    if !existe { return false }
    c.Rol = rolAlumno
    return true
}
\end{lstlisting}

Las operaciones de lectura usan \texttt{RLock}/\texttt{RUnlock} (múltiples lectores
simultáneos); las de escritura usan \texttt{Lock}/\texttt{Unlock} (exclusividad total).
\texttt{defer} garantiza que el mutex se libere siempre al salir de la función, incluso
si ocurre un panic.

\begin{lstlisting}[style=golang, caption={Control de sincronizacion}]
func (e *estadoGlobal) ActivarSync(p parametrosSync) {
    e.mu.Lock()
    defer e.mu.Unlock()
    e.sync.activo = true
    e.sync.parametros = p
}

func (e *estadoGlobal) DesactivarSync() {
    e.mu.Lock()
    defer e.mu.Unlock()
    e.sync.activo = false
    e.sync.parametros = parametrosSync{}
}

func (e *estadoGlobal) EstadoSync() (bool, parametrosSync) {
    e.mu.RLock()
    defer e.mu.RUnlock()
    return e.sync.activo, e.sync.parametros
}
\end{lstlisting}

El mecanismo de sincronización es un canal de difusión pasivo: el asistente o el
profesor escriben los parámetros una vez, y todos los alumnos los reciben la próxima
vez que llaman a \texttt{/solve}. No hay notificaciones push ni websockets; la app
consulta periódicamente y detecta el cambio por el campo \texttt{"sync"} de la
respuesta.

% -------- handlers.go -------------------------------------------
\subsection{\texttt{handlers.go} — Endpoints de la app móvil}

Este archivo implementa los cuatro endpoints con los que interactúa la aplicación
Unity: registro, cálculo y control de sincronización desde el rol de asistente.

\begin{lstlisting}[style=golang, caption={\texttt{handlers.go} — registro de alumno}]
package main

import (
    "net/http"
    "strconv"
    "github.com/gin-gonic/gin"
)

type registerRequest struct {
    Legajo string `json:"legajo"`
    Nombre string `json:"nombre"`
}

func registerHandler(c *gin.Context) {
    var req registerRequest
    if err := c.ShouldBindJSON(&req); err != nil {
        c.JSON(http.StatusBadRequest,
            gin.H{"error": "Cuerpo de solicitud invalido"})
        return
    }
    if req.Legajo == "" || req.Nombre == "" {
        c.JSON(http.StatusBadRequest,
            gin.H{"error": "Legajo y nombre son obligatorios"})
        return
    }
    estado.RegistrarCliente(c.ClientIP(), req.Legajo, req.Nombre)
    c.JSON(http.StatusOK, gin.H{"ok": true})
}
\end{lstlisting}

\texttt{ShouldBindJSON} deserializa el cuerpo de la petición en el struct e informa el
error si el JSON es malformado. La validación de campos vacíos es explícita porque
\texttt{ShouldBindJSON} no valida semántica, solo sintaxis. La IP del cliente la
provee Gin directamente con \texttt{c.ClientIP()}.

\begin{lstlisting}[style=golang, caption={\texttt{handlers.go} — calculo de la EDP}]
func solveHandler(c *gin.Context) {
    estado.ActualizarRequest(c.ClientIP())

    syncActivo, params := estado.EstadoSync()

    var L, T, cw float64
    var inicial, ecuacion string

    if syncActivo {
        L        = params.L
        T        = params.T
        cw       = params.C
        inicial  = params.Inicial
        ecuacion = params.Ecuacion
    } else {
        Ls, ok1 := c.GetQuery("L")
        Ts, ok2 := c.GetQuery("T")
        cs, ok3 := c.GetQuery("c")
        inicial, _  = c.GetQuery("inicial")
        ecuacion, _ = c.GetQuery("ecuacion")

        if !ok1 || !ok2 || !ok3 {
            c.JSON(http.StatusBadRequest,
                gin.H{"error": "Faltan parametros en la consulta"})
            return
        }

        var err1, err2, err3 error
        L,  err1 = strconv.ParseFloat(Ls, 64)
        T,  err2 = strconv.ParseFloat(Ts, 64)
        cw, err3 = strconv.ParseFloat(cs, 64)

        if err1 != nil || err2 != nil || err3 != nil {
            c.JSON(http.StatusBadRequest, gin.H{"error": "Parametros invalidos"})
            return
        }

        lMax := 10.0
        if ecuacion == "calor2d" { lMax = 20.0 }
        if L < 1 || L > lMax || T < 1 || T > 60 || cw < 0.1 || cw > 2 {
            c.JSON(http.StatusBadRequest,
                gin.H{"error": "Parametros fuera de rango"})
            return
        }
    }

    var result interface{}
    switch ecuacion {
    case "calor2d":
        result = ecuacion_calor_2d(L, T, cw, inicial)
    default:
        ecuacion = "onda2d"
        result = ecuacion_onda_2d(L, T, cw, inicial)
    }

    c.JSON(http.StatusOK, gin.H{
        "sync":     syncActivo,
        "ecuacion": ecuacion,
        "result":   result,
    })
}
\end{lstlisting}

Este handler implementa la lógica central del sistema. Lo primero que hace es consultar
el estado de sincronización. Si está activo, ignora completamente los query parameters
de la URL y usa los parámetros del asistente. Si no, los lee, los convierte de string a
float64 (\texttt{strconv.ParseFloat}) y valida rangos. El límite de \texttt{L} varía
según la ecuación: 10 m para onda 2D (donde la grilla crece como $n = 8L$) y 20 m para
calor 2D (donde crece como $n = 5L$), controlando el tamaño máximo de la grilla.

\begin{lstlisting}[style=golang, caption={\texttt{handlers.go} — sincronizacion desde asistente}]
type syncSetRequest struct {
    Ecuacion string  `json:"ecuacion"`
    L        float64 `json:"L"`
    T        float64 `json:"T"`
    C        float64 `json:"c"`
    Inicial  string  `json:"inicial"`
}

func syncSetHandler(c *gin.Context) {
    cliente := estado.BuscarCliente(c.ClientIP())
    if cliente == nil || cliente.Rol != rolAsistente {
        c.JSON(http.StatusForbidden, gin.H{"error": "Se requiere rol de asistente"})
        return
    }
    var req syncSetRequest
    if err := c.ShouldBindJSON(&req); err != nil {
        c.JSON(http.StatusBadRequest,
            gin.H{"error": "Cuerpo de solicitud invalido"})
        return
    }
    lMax := 10.0
    if req.Ecuacion == "calor2d" { lMax = 20.0 }
    if req.L < 1 || req.L > lMax || req.T < 1 || req.T > 60 ||
        req.C < 0.1 || req.C > 2 {
        c.JSON(http.StatusBadRequest, gin.H{"error": "Parametros fuera de rango"})
        return
    }
    estado.ActivarSync(parametrosSync{
        Ecuacion: req.Ecuacion,
        L: req.L, T: req.T, C: req.C, Inicial: req.Inicial,
    })
    c.JSON(http.StatusOK, gin.H{"ok": true})
}

func syncClearAsistenteHandler(c *gin.Context) {
    cliente := estado.BuscarCliente(c.ClientIP())
    if cliente == nil || cliente.Rol != rolAsistente {
        c.JSON(http.StatusForbidden, gin.H{"error": "Se requiere rol de asistente"})
        return
    }
    estado.DesactivarSync()
    c.JSON(http.StatusOK, gin.H{"ok": true})
}
\end{lstlisting}

Ambos handlers verifican el rol antes de cualquier otra operación. Si la IP no existe
en el mapa de clientes (dispositivo no registrado) o su rol no es asistente, se rechaza
con 403. La verificación de rol es la única diferencia con los endpoints equivalentes
de \texttt{admin.go}.

% -------- admin.go ----------------------------------------------
\subsection{\texttt{admin.go} — Endpoints del panel docente}

Este archivo implementa las rutas del grupo \texttt{/admin}. Su lógica es análoga a
\texttt{handlers.go} pero sin verificación de rol de asistente, ya que el acceso al
grupo \texttt{/admin} implica que el profesor ya fue autenticado por IP.

\begin{lstlisting}[style=golang, caption={\texttt{admin.go} — lista de clientes}]
package main

import (
    "net/http"
    "github.com/gin-gonic/gin"
)

type clienteResponse struct {
    IP            string `json:"ip"`
    Legajo        string `json:"legajo"`
    Nombre        string `json:"nombre"`
    Rol           string `json:"rol"`
    UltimoRequest string `json:"ultimoRequest"`
}

func clientesHandler(c *gin.Context) {
    clientes := estado.Clientes()
    lista := make([]clienteResponse, 0, len(clientes))
    for _, cl := range clientes {
        rolStr := "alumno"
        if cl.Rol == rolAsistente { rolStr = "asistente" }
        lista = append(lista, clienteResponse{
            IP:            cl.IP,
            Legajo:        cl.Legajo,
            Nombre:        cl.Nombre,
            Rol:           rolStr,
            UltimoRequest: cl.UltimoRequest.Format("15:04:05"),
        })
    }
    c.JSON(http.StatusOK, gin.H{"clientes": lista})
}
\end{lstlisting}

Se define un struct \texttt{clienteResponse} separado del \texttt{Cliente} interno.
Esta separación es deliberada: el struct interno usa \texttt{rol} (entero) y
\texttt{time.Time}, pero la API expone strings legibles. Si el struct interno cambia,
el contrato de la API no se ve afectado. \texttt{UltimoRequest.Format("15:04:05")}
formatea el timestamp como hora en formato \texttt{HH:MM:SS}.

\begin{lstlisting}[style=golang, caption={\texttt{admin.go} — promote, demote y sync}]
type promoverRequest struct {
    IP string `json:"ip"`
}

func promoverHandler(c *gin.Context) {
    var req promoverRequest
    if err := c.ShouldBindJSON(&req); err != nil {
        c.JSON(http.StatusBadRequest,
            gin.H{"error": "Cuerpo de solicitud invalido"})
        return
    }
    if !estado.PromoverAsistente(req.IP) {
        c.JSON(http.StatusNotFound, gin.H{"error": "Cliente no encontrado"})
        return
    }
    c.JSON(http.StatusOK, gin.H{"ok": true})
}

func demoverHandler(c *gin.Context) {
    var req promoverRequest
    if err := c.ShouldBindJSON(&req); err != nil {
        c.JSON(http.StatusBadRequest,
            gin.H{"error": "Cuerpo de solicitud invalido"})
        return
    }
    if !estado.DemoverAsistente(req.IP) {
        c.JSON(http.StatusNotFound, gin.H{"error": "Cliente no encontrado"})
        return
    }
    c.JSON(http.StatusOK, gin.H{"ok": true})
}

func syncClearHandler(c *gin.Context) {
    estado.DesactivarSync()
    c.JSON(http.StatusOK, gin.H{"ok": true})
}

func syncSetAdminHandler(c *gin.Context) {
    var req syncSetRequest
    if err := c.ShouldBindJSON(&req); err != nil {
        c.JSON(http.StatusBadRequest,
            gin.H{"error": "Cuerpo de solicitud invalido"})
        return
    }
    lMax := 10.0
    if req.Ecuacion == "calor2d" { lMax = 20.0 }
    if req.L < 1 || req.L > lMax || req.T < 1 || req.T > 60 ||
        req.C < 0.1 || req.C > 2 {
        c.JSON(http.StatusBadRequest, gin.H{"error": "Parametros fuera de rango"})
        return
    }
    estado.ActivarSync(parametrosSync{
        Ecuacion: req.Ecuacion,
        L: req.L, T: req.T, C: req.C, Inicial: req.Inicial,
    })
    c.JSON(http.StatusOK, gin.H{"ok": true})
}
\end{lstlisting}

\texttt{promoverHandler} y \texttt{demoverHandler} reutilizan el mismo struct
\texttt{promoverRequest} porque ambas operaciones solo necesitan la IP del cliente
a modificar. \texttt{syncClearHandler} no necesita validar nada: el profesor ya está
autenticado y desactivar una sincronización inexistente es una operación inocua.
\texttt{syncSetAdminHandler} es equivalente a \texttt{syncSetHandler} de la app móvil,
con la diferencia de que no verifica rol de asistente.

% ── SECCIÓN 5: TRABAJO FUTURO ───────────────────────────────
\newpage

\section{Trabajo Futuro}


El presente proyecto establece una base funcional para la simulación interactiva de
EDPs en realidad virtual, pero abre también varias líneas de trabajo que exceden su
alcance actual y que podrían ser abordadas en etapas posteriores.

\subsection{Escalabilidad de la plataforma}

La arquitectura actual está concebida para un entorno de aula con decenas de usuarios
simultáneos conectados a una red local. Escalar la plataforma a un uso más amplio
implicaría abordar al menos los siguientes aspectos:

\begin{itemize}
    \item \textbf{Persistencia de sesiones}: el estado de los clientes se mantiene en
    memoria volátil del servidor. Incorporar una base de datos permitiría sobrevivir
    reinicios y registrar el historial de simulaciones de cada usuario.

    \item \textbf{Cómputo distribuido}: para dominios grandes o parámetros que
    requieran grillas de alta resolución, los solvers podrían distribuirse entre
    múltiples núcleos o nodos mediante goroutines paralelas o una cola de tareas,
    reduciendo los tiempos de respuesta bajo carga concurrente.

    \item \textbf{Despliegue en la nube}: migrar el servidor a una instancia en la
    nube eliminaría la restricción de red local, permitiendo que estudiantes accedan
    a la plataforma desde cualquier ubicación. Esto requeriría además incorporar un
    sistema de autenticación robusto en reemplazo del esquema de IP actual.

    \item \textbf{Limitación de recursos por usuario}: implementar colas de trabajo y
    límites de tamaño de grilla por sesión para evitar que simulaciones costosas de
    un usuario degraden la experiencia del resto.
\end{itemize}

\subsection{Incorporación de nuevos solvers}

La arquitectura modular del servidor facilita la incorporación de nuevas ecuaciones
diferenciales. Algunas extensiones de interés pedagógico son:

\begin{itemize}
    \item \textbf{Ecuación de Laplace / Poisson}: resolución de problemas electrostáticos
    y de potencial en estado estacionario mediante métodos iterativos (Jacobi,
    Gauss-Seidel, SOR). Permitiría visualizar distribuciones de campo eléctrico en
    dominios con geometrías arbitrarias.

    \item \textbf{Ecuación de onda 1D}: versión simplificada de la membrana vibrante,
    ideal como paso introductorio antes de los casos 2D. Su baja carga computacional
    permitiría resoluciones muy finas y tiempos de simulación largos.

    \item \textbf{Navier-Stokes simplificado}: simulación de flujo de fluidos
    incompresible en 2D (ecuaciones de corriente-vorticidad). Abre la plataforma a
    cursos de mecánica de fluidos y dinámica de fluidos computacional.

    \item \textbf{Esquemas implícitos}: para la ecuación de calor, un esquema de
    Crank-Nicolson o totalmente implícito eliminaría la restricción de estabilidad
    sobre $\Delta t$, permitiendo simular escalas temporales mucho más largas sin
    incrementar el cómputo por paso.
\end{itemize}

\subsection{Mejoras en la experiencia de realidad virtual}

Más allá de los solvers, la experiencia inmersiva podría enriquecerse con:

\begin{itemize}
    \item Interacción gestual para modificar parámetros directamente dentro del entorno
    VR, sin salir de la aplicación.
    \item Representaciones volumétricas para problemas 3D, aprovechando la profundidad
    espacial que ofrece la realidad virtual.
    \item Modo colaborativo en tiempo real, donde múltiples usuarios en el entorno VR
    puedan modificar condiciones iniciales conjuntamente y observar el efecto de forma
    sincrónica.
\end{itemize}

\vfill

\end{document}
